\documentclass[a4paper,11pt]{article}
\usepackage{exam-style-842-unit}
\examinehead{842 信息化概论 · 第06单元练习册}{练习二 参考答案与解析}
\begin{document}

\unitanswerscover{842 信息化概论 · 单元练习册}{第06单元参考答案}{练习二（传感器、RFID 与传感器网络）的答案与解析；教材未展开的 RFID、WSN 系统细节按业界通行表述补充并在解析中注明来源；每道判断题均给出改正表述。}

\answerpart{一、单项选择题（10 题，每题 3 分，共 30 分）}

\q A。\begin{analysis}微电子技术是以集成电路为核心的电子技术，是在电子元器件微型化、低功耗过程中发展起来的，至今已经历了五代技术的突破与发展；计算机、数控机床、机器人、智能手机、家电设备、航空航天、军事装备等无一不使用微电子集成电路技术，因此微电子技术是传感技术、通信技术、计算机技术发展的重要基础（教材 5.1、5.1.1，第 65—66 页）。B、C、D 都是微电子技术支撑的对象，不是其核心。\end{analysis}

\q A。\begin{analysis}传感器（Sensor/Transducer），也称为换能器，是一种能感受到物理、化学、生物量信息的器件，并将感受到的信息按一定规律变换成电压信号、电流信号、脉冲信号等形式的输出信息，以满足信息的电子化传输、处理、存储、显示、记录和控制等需要（教材 5.2.2，第 74 页）。国家标准 GB7665-87 的定义是：能感受规定的被测量并按照一定的规律转换成可用信号的器件或装置，通常由敏感元件和转换元件组成。B 执行器是执行机构；C 控制器是控制核心；D 显示器是输出设备。\end{analysis}

\q A。\begin{analysis}敏感元件（预变换器）是能直接感受被测量（一般为非电量）并输出与被测量成确定关系的其他物理量的元件。例如，压力传感器中的膜片就是敏感元件，它首先将压力转换为位移，然后再将位移转换为电量；对于不能直接变换为电量的传感器必须进行预变换（教材 5.2.2，第 74—75 页）。B 转换元件负责把敏感元件输出量变为电量；C 转换电路把电参量变为电压、电流或频率；D 辅助电源供电。\end{analysis}

\q B。\begin{analysis}灵敏度指在稳态时输出量的变化量与输入量变化量的比值，即单位输入量引起的输出量大小，这个比值越大，传感器就越灵敏（教材 5.2.4，第 77—78 页）。A 量程是允许误差限范围内被测量的测量范围；C 线性度是校准曲线与拟合直线的偏离程度；D 迟滞是正、反向特性曲线的不一致程度。\end{analysis}

\q C。\begin{analysis}迟滞现象是传感器正向特性曲线（输入量增大）和反向特性曲线（输入量减小）的不一致程度，即由小到大的测量与由大到小的测量过程总是存在差异性（教材 5.2.4，第 79 页）。A 不重复性是在相同条件下按同一方向做全量程多次重复测量时输出曲线的不一致程度；B 线性度是校准曲线与拟合直线的偏离程度；D 零漂是输入不变或为零时输出随时间变化的特性。\end{analysis}

\q B。\begin{analysis}精确度也称为精度，是线性度、不重复性及迟滞 3 项指标的综合指数，反映了系统误差和随机误差的综合指标（教材 5.2.4，第 80 页）。A 灵敏度是输出变化量与输入变化量之比；C 分辨率不是教材所列 8 个静态特性指标之一；D 量程是测量范围。\end{analysis}

\q B。\begin{analysis}智能传感器是由传统的传感器和微处理器（或微计算机）相结合而构成的，微处理器是智能传感器的核心，不但可以对传感器的测量数据进行计算、存储、处理，还可以通过反馈回路对传感器进行调节（教材 5.2.5，第 81 页）。A、C、D 是传统传感器的组成部件，不能承担判断、信息处理与记忆等智能功能。\end{analysis}

\q A。\begin{analysis}RFID 系统由电子标签、读写器、天线和后端信息系统组成。电子标签附着在物品上，存储规范而具有互用性的信息，是物品的“身份证”；读写器通过天线发射射频信号，读取或写入标签信息；后端信息系统对采集到的数据进行处理和应用（RFID 系统组成按业界通行表述补充；教材 7.1.2 第 115 页指出 RFID 标签上存储着规范而具有互用性的信息、RFID 由标签和阅读器组成）。B 读写器负责读取或写入；C 天线负责射频信号的发射与接收；D 后端信息系统负责数据处理。\end{analysis}

\q A。\begin{analysis}无线传感器网络节点通常由电池供电，部署后往往难以更换电池，因此能量是节点最突出的资源约束，节能是网络协议与算法设计的核心目标（WSN 节点组成与能量约束按业界通行表述补充；教材 5.2.6 第 84 页指出无线智能传感器与通信技术融合将形成泛在无线传感器网络，是物联网发展的关键）。B、C、D 是节点的其他资源，但都不是最突出的约束。\end{analysis}

\q A。\begin{analysis}1947 年，美国贝尔实验室的物理学家肖克利、巴丁和布拉顿三人合作发明研制了晶体管——一种由 3 个支点构成的半导体固体元件，晶体管当时被称为“三条腿的魔术师”，三人因此共同获得 1956 年诺贝尔物理学奖（教材 5.1.1，第 66 页）。B 基尔比和诺伊斯分别发明了锗集成电路和硅集成电路；C 香农、维纳、贝塔朗菲分别是信息论、控制论、系统论的代表人物；D 盖茨、乔布斯、艾伦与晶体管发明无关。\end{analysis}

\answerpart{二、多项选择题（4 题，每题 4 分，共 16 分）}

\q ABC。\begin{analysis}传感器一般由敏感元件与转换元件两个基本元件组成；在完成非电量到电量的变换过程中，实际应用还包括转换电路和辅助电源：转换电路把转换元件输出的电参量转换成电压、电流或频率量，辅助电源为需要电源才能工作的转换电路和转换元件提供正常工作电源（教材 5.2.2，第 74—75 页）。A、B、C 正确。D 错误：若转换元件输出的已经是电压、电流或频率等电参量，就不用转换电路，因此转换电路并非任何传感器都必须配置。\end{analysis}

\q ABC。\begin{analysis}传感器一般按 4 种方法分类：按被测物理量分为位移量、力量、运动量、热学量、光学量、气体量等传感器；按工作原理分为电阻式、电容式、电感式、压电式、霍尔式、光电式、光栅式、热电式等传感器；按输出信号性质分为开关型、数字型和模拟型传感器；按照能量传递方式分为有源传感器与无源传感器两大类（教材 5.2.3，第 75—76 页）。A、B、C 正确。D 错误：按生产厂家规模分类不是教材给出的分类方法。\end{analysis}

\q ABC。\begin{analysis}与条形码相比，RFID 是非接触式识别，不需要人工逐件近距离扫描；标签信息可重复读写、可加密；读写器可同时批量读取多个标签，作用距离从几厘米到几十米，且抗污损能力强，适合畜牧、物流等恶劣环境（RFID 与条码比较按业界通行表述补充；教材 7.1.2 第 115 页指出 RFID 标签信息可通过无线数据通信网自动采集到中央信息系统，实现物品识别）。A、B、C 正确。D 错误：条形码印刷成本低，RFID 标签成本通常高于条形码，成本低是条形码的优势而非 RFID 的优势。\end{analysis}

\q ABC。\begin{analysis}无线传感器网络由大量传感器节点构成，节点通过自组织方式组网，数据以多跳方式路由传输到汇聚节点；网络以数据为中心，关注的是监测区域内的数据而非节点地址；节点通常电池供电、能量受限，这是 WSN 区别于传统网络的重要特点（按业界通行表述补充；教材 5.2.6 第 84 页提出泛在无线传感器网络是物联网发展的关键）。A、B、C 正确。D 错误：节点能量通常不受限制的说法与 WSN 实际相反。\end{analysis}

\answerpart{三、判断题（6 题，每题 2 分，共 12 分）}

\q √。\begin{analysis}教材 5.1、5.1.1 指出：微电子技术是以集成电路为核心的电子技术，是在电子元器件微型化、低功耗过程中发展起来的，至今已经历了五代技术的突破与发展——电子管技术、晶体管技术、集成电路、集成电路的规模化发展、微机电集成电路（教材 5.1.1，第 65—66 页）。\end{analysis}

\q √。\begin{analysis}传感器一般由敏感元件与转换元件两个基本元件组成；在完成非电量到电量的变换过程中，实际应用还包括转换电路和辅助电源（教材 5.2.2，第 74—75 页）。该项表述与教材一致。\end{analysis}

\q ×。改正：传感器在相同条件下按同一方向做全量程多次重复测量时输出曲线的不一致程度称为不重复性；迟滞是指传感器正向特性曲线（输入量增大）与反向特性曲线（输入量减小）的不一致程度。\begin{analysis}教材 5.2.4 把不重复性与迟滞列为两个不同的静态特性指标：不重复性强调“同一方向、多次重复测量”的一致性，迟滞强调“正程与反程”的差异（教材 5.2.4，第 79 页）。混淆二者是本单元最常见的失分点。\end{analysis}

\q √。\begin{analysis}智能传感器是由传统的传感器和微处理器（或微计算机）相结合而构成的，微处理器是智能传感器的核心，可以对测量数据进行计算、存储、处理，并通过反馈回路对传感器进行调节（教材 5.2.5，第 81 页）。\end{analysis}

\q ×。改正：RFID 是非接触式自动识别技术，通过射频信号非接触地读取或写入标签信息，作用距离可从几厘米到几十米，不需要贴近读写器。\begin{analysis}RFID 利用射频信号的空间耦合（电磁感应或电磁传播）实现非接触信息传递，这是它区别于条形码等接触式识别的根本特点（RFID 非接触原理按业界通行表述补充；教材 7.1.1 第 113 页把 RFID 列为物联网定义中的信息传感设备，7.1.2 第 115 页指出其信息可自动采集到中央信息系统）。\end{analysis}

\q ×。改正：无线传感器网络节点通常由电池供电、能量受限，节能是网络设计的核心问题之一，不能长期持续高速传输数据；通常需要采用休眠调度、分簇与数据融合、短距离多跳等节能措施。\begin{analysis}WSN 节点体积小、成本低，通常依靠电池供电且部署后难以更换电池，能量约束直接决定网络寿命（WSN 能量约束按业界通行表述补充；教材 5.2.6 第 84 页把网络化智能传感器列为发展方向，其关键是智能传感、野外供能和自动组网等技术）。\end{analysis}

\answerpart{四、填空题（6 空，每空 1 分，共 6 分）}

\q 敏感元件。\begin{analysis}能直接感受被测量（一般为非电量）并输出与被测量成确定关系的其他物理量的元件称为敏感元件，又称预变换器（教材 5.2.2，第 74—75 页）。\end{analysis}

\q 输入量变化量。\begin{analysis}灵敏度指在稳态时输出量的变化量与输入量变化量的比值，即单位输入量引起的输出量大小（教材 5.2.4，第 77—78 页）。\end{analysis}

\q 精确度（精度）。\begin{analysis}精确度也称为精度，是线性度、不重复性及迟滞 3 项指标的综合指数（教材 5.2.4，第 80 页）。\end{analysis}

\q 读写器（阅读器）。\begin{analysis}RFID 系统由电子标签、读写器、天线和后端信息系统组成；教材 7.1.2 第 115 页简称为“标签和阅读器”。读写器负责通过天线发射射频信号并读取或写入标签信息。\end{analysis}

\q WSN。\begin{analysis}无线传感器网络的英文缩写是 WSN（Wireless Sensor Network）；教材 5.2.6 第 84 页提出了泛在无线传感器网络的发展方向，并指出其关键是智能传感、野外供能和自动组网等技术。\end{analysis}

\q 微机电系统。\begin{analysis}MEMS（Micro Electro Mechanical Systems）即微机电系统，从广义上讲是指集微型传感器、微型执行器、信号处理器和控制电路、接口电路、通信系统以及电源于一体的微型机电系统（教材 5.3.1，第 85 页）。\end{analysis}

\answerpart{五、简答题（4 题，每题 6 分，共 24 分）}

\q 参考答案要点：
（1）定义（2 分）：传感器（Sensor/Transducer）又称换能器，是一种能感受到物理、化学、生物量信息的器件，并将感受到的信息按一定规律变换成电压信号、电流信号、脉冲信号等形式的输出信息，以满足信息的电子化传输、处理、存储、显示、记录和控制等需要。国家标准 GB7665-87 的定义为：能感受规定的被测量并按照一定的规律转换成可用信号的器件或装置，通常由敏感元件和转换元件组成。通俗地说，传感器是能把外界非电量信息转换成电量输出的器件。（教材 5.2.2，第 74 页）
（2）组成及作用（3 分）：
敏感元件（预变换器）：直接感受被测量并输出与被测量成确定关系的其他物理量，完成预变换，如压力传感器中的膜片先把压力转换为位移；
转换元件（传感元件）：把敏感元件的输出量转换成电量后再输出，如把位移量变换为电容量、电阻量或电感量；
转换电路：把转换元件输出的电参量转换成电压、电流或频率量，如电阻应变片传感器采用的电桥测量电路；若转换元件输出已是上述电参量则不用此电路；
辅助电源：为需要电源才能工作的转换电路和转换元件提供正常工作电源，通常采用直流电源。
（3）意义（1 分）：传感器是人类感觉器官的延伸或拓展，故又称为电子器官；其作用可归纳为对被测信号产生敏感并把它转换出来、在信号提取的同时把它转换成所需要的信号（主要是完成非电量到电量的转换）。（教材 5.2.3，第 75 页）

\q 参考答案要点（教材 5.2.4，第 77—80 页）：
（1）量程：在允许误差限范围内被测量值的测量范围。
（2）灵敏度：稳态时输出量的变化量与输入量变化量的比值，即单位输入量引起的输出量大小，比值越大传感器越灵敏。
（3）线性度：传感器的校准曲线与选定的拟合直线的偏离程度，又称非线性误差。
（4）不重复性：在相同条件下传感器的输入量按同一方向做全量程多次重复测量，输出曲线的不一致程度；重复测量的复合度越高，测量一致性越好。
（5）迟滞：传感器正向特性曲线（输入量增大）和反向特性曲线（输入量减小）的不一致程度。
（6）精确度：又称精度，是线性度、不重复性及迟滞 3 项指标的综合指数，反映系统误差和随机误差的综合指标。
（7）零点时间漂移（零漂）：在恒定温度环境中，当输入信号不变或为零时，输出信号随时间变化的特性。
（8）零点温度漂移（温漂）：当输入信号不变或为零时，传感器的输出信号随温度变化的特性。
补充：衡量传感器静态特性好坏主要用上述 8 个指标；一个高精度传感器必须有良好的静态特性和动态特性，才能完成信号无失真的转换与测量。常见失分：把不重复性与迟滞混为一谈；漏掉零漂、温漂两项。

\q 参考答案要点：
（1）概念（1 分）：RFID 即射频识别技术，通过射频信号非接触地自动识别目标对象并获取相关数据，是物联网中让物品“开口说话”的关键技术。（教材 7.1.2，第 115 页）
（2）系统组成（2 分，按业界通行表述补充，教材 7.1.2 只指出 RFID 由标签和阅读器组成）：电子标签，附着在物品上、存储规范而具有互用性的信息，分无源标签和有源标签；读写器，通过天线发射射频信号，读取或写入标签信息，并将数据送后端系统；天线，负责射频信号的发射与接收；后端信息系统，对采集到的标签数据进行存储、处理和应用。
（3）工作原理（1 分）：读写器经天线发射射频信号形成电磁场；标签进入磁场后，无源标签通过电磁感应获取能量并被激活，有源标签则依靠自带电源工作；标签把存储的信息调制后反射（回传）给读写器；读写器接收、解调后送后端信息系统处理，完成非接触识别。
（4）与条形码比较的优点（2 分）：非接触识别，不需要人工逐件近距离扫描；信息可重复读写、可加密；可批量读取多个标签；作用距离远、抗污损能力强，适合畜牧、物流等环境。条形码的优点是印刷简单、成本低，信息一般不可改写。
（5）常见失分：只答“标签＋读写器”而漏掉天线和后端信息系统；把 RFID 说成接触式识别。

\q 参考答案要点：
（1）概念（1 分）：无线传感器网络（WSN）是由大量传感器节点以自组织、多跳方式构成的无线网络，节点协作感知、采集和处理覆盖区域内被监测对象的信息，并把数据传送到汇聚节点（按业界通行表述补充；教材 5.2.6 第 84 页提出泛在无线传感器网络是物联网发展的关键）。
（2）节点组成（2 分）：传感器模块，负责感知和采集被监测对象的信息；处理器模块，负责数据处理、存储与协议运行；无线通信模块，负责与其他节点或汇聚节点通信；能量供应模块，通常为电池，为节点供电。
（3）主要特点（2 分）：自组织组网，无需固定基础设施；多跳路由传输，节点既是终端又充当路由器；动态拓扑，节点加入、退出或失效时网络可自愈；以数据为中心，关注监测数据而非节点地址；大规模密集部署；能量受限。
（4）节能思路（1 分，答出 3 条即可）：采用休眠/唤醒的占空比调度，减少空闲监听；分簇路由与数据融合，减少冗余传输；采用短距离多跳和低功耗 MAC 协议（如 IEEE 802.15.4/ZigBee 类技术）；降低采样与上报频率、分级上报；使用能量收集或太阳能补充供电。
（5）常见失分：把节点组成写成“传感器＋电池”两项；只说“省电”而不给具体机制。

\answerpart{六、综合题（1 题，共 12 分）}

\q 参考答案要点：
（1）数据采集链条（4 分，每环节 1 分）：被测量是牛舍空气中的氨气浓度（化学量）；敏感元件是气敏材料/气敏元件，氨气浓度变化使其电阻（电导率）发生确定关系的变化，完成“非电量到易于转换的非电量或电参量”的预变换；转换元件把该电阻变化转换成电压或电流信号；转换电路（电桥、放大电路等）把电参量转换成便于传输、处理的电压、电流或频率量；最后由采集接口输出电量信号并送往无线通信模块。若不能一次直接变换为电量，则先由敏感元件预变换，再由转换元件变换为电量，这正是教材 5.2.2 所述“并非所有的非电量参数都能一次直接变换为电量”的含义。（教材 5.2.2，第 74—75 页）
（2）RFID 系统组成与无源标签获能（4 分）：
组成——电子标签（电子耳标，存储牛号等规范信息，附着在奶牛身上）；读写器（安装在挤奶厅入口，读取耳标信息）；天线（发射和接收射频信号）；后端信息系统（云端管理平台，存储处理奶牛身份、产奶、健康数据）。
无源标签获能——读写器天线发射射频信号，在周围形成电磁场；无源电子耳标进入场区后，通过电磁感应（近场）或电磁传播（远场）从射频场中获取能量，激活芯片；芯片把存储信息调制后反射回读写器，读写器解调后送后端系统，完成非接触识别。整个过程耳标不需要自带电池，因而体积小、寿命长、成本低。（教材 7.1.2 第 115 页指出 RFID 标签信息可由无线数据通信网自动采集到中央信息系统；系统组成与无源获能机理按业界通行表述补充）
（3）电子耳标优势、MEMS 特点与节能措施（3 分）：
优势（1 分）——非接触、远距离读取，奶牛不必逐头靠近扫描；可重复读写，便于更新免疫、产奶记录；可批量读取，提高挤奶厅通行效率；抗污损、抗潮湿，适合牛舍环境；可加密，防止标签信息被伪造。
MEMS 特点（1 分）——微型化：器件体积小、质量轻、功耗低、响应时间短，具有微米、纳米精度的加工和毫米级体积；集成化：可把不同功能的多个传感器或执行器集成于一体，形成微传感器阵列；低成本：用硅微加工工艺在一片硅片上可同时制造成百上千个器件，适合批量生产。牛舍中可用 MEMS 加速度计监测奶牛活动量、用 MEMS 压力传感器监测管道压力。（教材 5.3.1，第 85—86 页）
节能措施（1 分，答出 3 条即可）——节点采用休眠/唤醒的占空比调度，无事件时休眠；分簇路由并由簇头做数据融合，减少冗余上报；采用 ZigBee 等短距离、低功耗通信与多跳传输，降低发射功率；按阈值触发上报（如氨气浓度超限才上报），减少通信次数；使用能量收集或太阳能电池补充供电。
（4）结论（1 分）：该奶牛场是传感器数据采集、RFID 个体识别、MEMS 与无线传感器网络相结合的典型农业物联网应用；答题时把“传感器组成—RFID 系统组成—MEMS 特点—WSN 节点与节能”几条线串起来，即可覆盖 2023 年论述题“农业物联网的关键技术”的采分点。（教材 5.2、5.2.2、5.2.5、5.3.1，第 73—86 页；7.1.2，第 115 页）

\end{document}
